\documentclass[journal]{IEEEtran}
\pdfminorversion=4
\usepackage{amsmath,amssymb,amsthm,bm,graphicx,booktabs,array,cite,url}
\usepackage[hidelinks]{hyperref}
\newtheorem{theorem}{Theorem}
\newtheorem{lemma}{Lemma}
\newtheorem{proposition}{Proposition}
\newtheorem{corollary}{Corollary}
\theoremstyle{definition}\newtheorem{definition}{Definition}
\theoremstyle{remark}\newtheorem{remark}{Remark}
\newcommand{\R}{\mathbb R}\newcommand{\E}{\mathbb E}
\newcommand{\K}{\mathcal K}\newcommand{\PiH}{\Pi_H}
\newcommand{\pos}[1]{[#1]_+}\newcommand{\norm}[1]{\left\lVert#1\right\rVert}
\newcommand{\ind}{\mathbf 1}\newcommand{\dd}{\mathrm d}
\newcommand{\clip}{\operatorname{clip}}\newcommand{\diag}{\operatorname{diag}}
\newcommand{\argmin}{\operatorname*{arg\,min}}
\hypersetup{pdftitle={Sharp Information Limits for Deterministic Two-Task Self-Tests Under Rate-Limited Healthy Drift},pdfauthor={},pdfkeywords={Planned self-test, deterministic scheduling, minimax, healthy drift}}
\begin{document}
\title{Sharp Information Limits for Deterministic Two-Task Self-Tests Under Rate-Limited Healthy Drift}
\author{Author details to be supplied}
\maketitle
\begin{abstract}
We study planned two-task self-tests with deterministic observation schedules, a specified onset and additive fault-growth direction, and known Gaussian readout covariance. Both hypotheses admit the same amplitude- and rate-constrained class of healthy paths; their least favorable paths may differ. Task switching reverses the healthy response while preserving the physical fault direction, but excludes transition readings and enlarges the reachable healthy increment. For a linearly growing fault and an unbounded healthy level, we determine the exact coefficient of the optimal five-halves-power information deficit relative to a full-calendar oracle. A uniform converse is attained by symmetric task blocks with square-root-scale holds, yielding a positive second-order penalty in the planned detection horizon. With a fixed finite healthy bound, the optimal deficit equals an explicit quadratic polynomial up to a bounded remainder, and staying loses at most a bounded amount of information. The constant-amplitude case gives a marginal holding-time rule. A closed-loop residual certificate reduces admission-tolerance and calendar selection to a finite optimization of protected power. In a slow-ramp, single-channel rigid-mechanism example, a 75-second symmetric experiment has protected power above 0.999 at level 0.001 while staying has zero minimax separation. Comparisons with fixed holds identify the intermediate-horizon benefit of growing holds. Ideal and deterministic-error-protected powers are reported together.
\end{abstract}
\begin{IEEEkeywords}
Fault diagnosis, experiment design, minimax testing, rate-limited uncertainty, switching cost, robot control.
\end{IEEEkeywords}
\section{Introduction}
\IEEEPARstart{P}{lanned} self-tests can deliberately move a controlled mechanism between configurations to discriminate a specified fault alternative from admissible healthy drift. This paper studies that design problem over a fixed observation horizon. The task calendar and fault onset are specified before data are observed; the objective is fixed-horizon minimax detection power, with optional preassigned multiple-test error allocation. The design is useful when a dominant healthy response reverses between two tasks while an additive fault response does not. A transition consumes observation time and permits the healthy parameter to change, so scheduling changes the uncertainty set as well as the sampling budget.

We determine the optimal information deficit over all allowed deterministic calendars. For linear fault growth with an unbounded healthy level, the deficit is asymptotic to a known constant times $H^{5/2}$. A schedule-uniform converse matches a symmetric half-window--full-window--half-window construction. For a fixed finite healthy interval, the deficit instead equals an explicit quadratic polynomial up to $O(1)$, and switching has only a bounded advantage over staying. The change of regime follows from the same physical-time drift constraint.

The system-design consequence is quantitative: transition time $\tau_m$ and the difference-path drift rate $R$ enter the second-order planned-time penalty through $\sqrt{\tau_m R/\kappa}$. A valid controller settling bound therefore changes a diagnostic cost with an explicit coefficient. Exact finite-history projections are used when only a few task blocks fit.

\subsection{Relation to previous work}
Experiment choice for hypothesis discrimination has classical sequential formulations due to Chernoff \cite{chernoff} and modern adaptive formulations in active hypothesis testing \cite{naghshvar}. Atkinson and Fedorov develop discriminating regression designs \cite{atkinson}. The predetermined calendar fixes the order of experiment choice and worst healthy-path choice; the sharp constant is specific to that order.

Minimax detection of uncertain signals in Gaussian noise \cite{burnashev} and convex Gaussian testing \cite{gjn} supply the statistical foundation. Likelihood profiling and nuisance rejection are established tools \cite{basseville}. The closest-pair test is an established tool; the new calculations concern the pathwise projection and its optimal dependence on a calendar with transition exclusions.

Robust detection and isolation of abrupt and incipient faults in uncertain nonlinear systems is an established control topic \cite{zhang2002}. Guaranteed active fault detection uses separating inputs for bounded uncertain models \cite{nik98,campbell,scott}; invariant residual sets and support-function duality also lead to active fault-isolation designs \cite{blanchini}. Heirung and Mesbah review this input-design perspective \cite{heirung}. Our noise-free limit is set separation in a two-task experiment. Joint stochastic control and diagnosis \cite{simandl} and the dependence of feedback benefits on the chosen criterion \cite{ashari} further motivate stating the controller and measurement interface separately.

Trend-free run orders with level-change costs \cite{coster}, autoregressive-error run orders \cite{cheng}, and correlated-error sampling designs \cite{bickel1,bickel2} are direct design antecedents. Radiometric timing also balances drift and dead time \cite{schieder,ossenkopf}. Here a transition changes the feasible increment of an unknown mean path, rather than only a counting cost or a prescribed random-error covariance. Independent optimization of each holding window would give a different adversary.

Sensor scheduling for dynamical estimation \cite{vitus,zhao}, observation control and switching costs in quickest change detection \cite{banerjee,lautay,lubenia}, and nonstationary information-growth analysis \cite{liang,hou} address related resource and time constraints. The present result is the sharp second-order deficit under a common rate-limited healthy-path class.

\subsection{Contributions and organization}
The three main results give constant-amplitude rates, sharp growth deficits, and finite-box switching advantages. An admission-design proposition couples a verified settling envelope with the protected statistical score and reduces the tolerance choice to finite calendar comparisons. In the slow-ramp example, the original declared readout standard deviation of $0.01$ N m already gives protected power above $0.999$ at 75 s under symmetric switching, while staying has zero information until about 110 s. Fixed-hold comparisons quantify the extra benefit of the growing-hold rule. A faster-ramp example then shows why a small residual norm alone does not guarantee a useful protected test. Physical calibration and information-interface assumptions are stated in Section~\ref{sec:scope}.

\section{Diagnostic Experiment}
\subsection{Time, schedules, and uncertainty}
The time-slot duration is $\delta>0$. A task transition excludes $k\ge1$ slots from the admissible observation model. Its duration is $\tau_m=k\delta$, whereas the interval between the last observation before an immediate transition and the first observation after it is $(k+1)\delta$. These two times are distinct.

\begin{definition}[Admissible deterministic calendar]\label{def:calendar}
After $n_0$ consecutive healthy-prefix readings at task $+$, a calendar $\pi\in\PiH$ specifies $H$ post-onset slots. At either of two tasks, each holding slot supplies one reading. A change to the other task uses at least $k$ consecutive slots with no admissible reading. The calendar, including the transition durations, is fixed before observing the data. Consecutive moves and empty round trips are allowed. A terminal move without a later reading is dominated by observing at the current task. The schedule \emph{stay} retains all prefix and post-onset observations at the initial task. A maximal contiguous sequence of readings at one task is a run.
\end{definition}

A physical observation at a task with sign $\epsilon_i\in\{+1,-1\}$ has the ideal form $\widetilde Y_i=\epsilon_i\beta^h(t_i)+a_i^h+\xi_i$. Multiplication by $\epsilon_i$ gives
\begin{equation}
 Y_i=\beta^h(t_i)+s_i^h+\widetilde\xi_i,
 \quad\widetilde\xi_i\stackrel{\rm iid}{\sim}N(0,\sigma^2).
 \label{eq:obs}
\end{equation}
The two hypotheses use independent choices $\beta^0,\beta^1$ from the same admissible class; equality of their realized paths is not imposed. Under $H_0$, $s^0=0$. Under $H_1$, prefix entries are zero and post-onset entries are $s_i^1=\epsilon_i a_{j_i}$, where $j_i$ is the physical slot index. In either hypothesis,
\begin{equation}
 |\beta^h(t)|\le b,\qquad |\dot\beta^h(t)|\le r_d
 \quad\text{a.e.}\label{eq:health}
\end{equation}
Initial healthy levels are unknown and free inside their bounds. The prefix supplies zero fault target, not a zero healthy-state boundary condition. The radii are upper envelopes: underestimating $b$ or $r_d$ may invalidate uniform size control, whereas enlarging a valid healthy class can only reduce its information. The independence and covariance assumptions are separate from these mean-path bounds.

Throughout, the one-sided bounds and their difference-class counterparts are
\begin{equation}
 B=2b,\quad R=2r_d,\quad r=R\delta,
 \qquad a_j=\rho_0+\eta j,\quad\eta=\kappa\delta.\label{eq:units}
\end{equation}
Thus $r$ and $\eta$ are changes per slot; $R$ and $\kappa$ are changes per unit physical time. We assume $\rho_0\ge0$ and, in the progressive results, $\eta>0$. Direction, onset, and growth are specified design inputs. The constant-amplitude case is stated separately.

\begin{lemma}[Exact calendar reduction]\label{lem:trace}
At the strictly increasing slot indices $j_1,\ldots,j_N$ including the prefix, the difference of the two symmetric healthy classes has the exact trace
\begin{equation}
 \K_\pi(B,r)=\{z: |z_i|\le B,\ |z_{i+1}-z_i|\le r(j_{i+1}-j_i)\}.
 \label{eq:K}
\end{equation}
\end{lemma}
\begin{proof}
Integration gives the increment inequalities. Conversely, piecewise-linear interpolation of admissible nodes has slope at most $R$ in physical time and stays inside the interval. Every difference path is realizable by $z/2$ and $-z/2$.\end{proof}
The optimal observed node vector will be unique, but its continuous extension through unobserved intervals need not be. No intermediate samples are introduced by the interpolation. Curvature constraints require velocity reachability; bounds on second differences alone are only a node-level relaxation.

\subsection{Information and the two oracles}
Let $s$ denote the stacked specified fault target, including its zero prefix. Define
\begin{equation}
 G_H^B(\pi)=\frac1{2\sigma^2}\min_{z\in\K_\pi(B,r)}\norm{s-z}^2.\label{eq:G}
\end{equation}
For $B=\infty$ the amplitude constraint is omitted. The closed finite-dimensional convex projection is still attained. The full-calendar and observed-calendar oracle informations are
\begin{equation}
 O_H=\frac1{2\sigma^2}\sum_{j=1}^Ha_j^2,
 \quad O_H^{\rm obs}(\pi)=\frac1{2\sigma^2}\sum_{j\in I_\pi}a_j^2.
\end{equation}
Consequently,
\begin{equation}
0\le G_H^B(\pi)\le O_H^{\rm obs}(\pi)\le O_H,
\quad D_H^{B,*}=O_H-\max_{\pi\in\PiH}G_H^B(\pi).\label{eq:D}
\end{equation}
Here $D_H^B(\pi)=O_H-G_H^B(\pi)$ denotes the deficit of a specific calendar. The finite calendar maximum exists. In particular,
\begin{equation}
 O_H=\frac{H\rho_0^2+\rho_0\eta H(H+1)+\eta^2H(H+1)(2H+1)/6}{2\sigma^2}.
 \label{eq:oracle}
\end{equation}

\begin{proposition}[Fixed-calendar testing]\label{prop:testing}
For a fixed calendar and specified target $s$, the best level-$\alpha$ worst-case power is
\begin{equation}
 \mathcal B_\pi(\alpha)=\Phi\bigl(\sqrt{2G_H^B(\pi)}-z_{1-\alpha}\bigr).\label{eq:power}
\end{equation}
\end{proposition}
\begin{proof}
If $s$ lies in the difference class, the mean sets share a law. Any level-$\alpha$ test has worst power at most $\alpha$, and constant randomized rejection attains it. Otherwise, let $z^*$ be the projection onto the difference class. The closest means are $\mu_0^*=z^*/2$ and $\mu_1^*=s-z^*/2$. The projection variational inequality makes $s-z^*$ a separating normal for both full mean classes. Its Gaussian threshold test is uniformly level $\alpha$ and attains the stated worst power at this pair. Neyman--Pearson for that pair gives the reverse inequality; this is the convex Gaussian testing reduction of \cite{gjn}, within the minimax uncertain-signal tradition \cite{burnashev}.\end{proof}
This is an exact minimax statement for the specified shift family. The distance is computed by profiling, but an arbitrary generalized likelihood-ratio statistic need not have the distribution in \eqref{eq:power}.

\subsection{Fault parity, calibration, and growth uncertainty}
For a fixed calendar and a pure ramp, write $s_{\min}=\kappa_{\min}u$ with $\kappa_{\min}>0$ and let $v=s_{\min}-z^*$. The level-$\alpha$ score is
\begin{equation}
 v^TY>\tfrac12v^Tz^*+\sigma\|v\|z_{1-\alpha}.
 \label{eq:implement_test}
\end{equation}
Its null class does not depend on $\kappa$. Since $0\in\K_\pi$, projection gives $v^Tz^*\ge0$, hence
$\kappa_{\min}v^Tu=\|v\|^2+v^Tz^*\ge\|v\|^2$.
The worst-case power of this \emph{fixed} test is therefore nondecreasing for $\kappa\ge\kappa_{\min}$ with the same healthy class and covariance. A nonzero fixed intercept or other template change requires checking $v^Tu\ge0$ explicitly. This statement does not equate a mistuned calendar's information with the sharp optimal coefficient.

\paragraph{Healthy prefix.}
The prefix is placed at task $+$ to fix the initial task and make its time cost explicit. A two-task prefix would consume an additional transition and could improve finite-horizon separation or calibration. It is a different initial calendar; the $H\to\infty$ leading coefficient is unchanged by a fixed amount of prefix work. A constant transformed offset is observable at either task, whereas learning unknown task-specific columns requires calibration data that the present model assumes available.

\section{Closed-Loop Implementation Interface}\label{sec:control}
The following rigid-mechanism example supplies a physical-time transition budget and a deterministic residual-error bound. Its diagnosis interface is explicitly restricted to selected torque-residual readings. Model-based robot residuals have established momentum-based and collision-detection formulations \cite{deluca,haddadin}; the result developed here concerns when a certified static residual may be used by a predetermined self-test.

\subsection{Dynamics, sensing, and the test alternative}
Consider the numerical planar revolute chain with terminal mass $m=\bar m+\beta_m$. A load-side additive bias acts along a specified joint direction $e_f$:
\begin{equation}
 \begin{split}
 M(q,m)\ddot q+c(q,\dot q,m)+g(q,m)\\
 +\dot m A(q)\dot q=\tau_u-e_fa(t),\qquad A=J_p^TJ_p.
 \end{split}\label{eq:plant}
\end{equation}
The $\dot m$ term specifies the time-dependent-Lagrangian convention. The chain has six coupled angular coordinates in one plane. It is a synthetic mechanism, rather than a model of a spatial industrial arm. Its reflected rotor inertia is $0.5$ kg m$^2$ per joint; for interpretation, $N^2J_m=0.5$ at reduction $N=100$ corresponds to $J_m=5\times10^{-5}$ kg m$^2$. No gearbox friction is included in this benchmark.

A concrete realization of the specified alternative is a scheduled torque-loading test: an output-side loading fixture applies an increasing, known-direction joint torque while the mechanism visits the two poses. It is a reproducible fault surrogate for an additive load-side bias. The linear law is a chosen test alternative, not a general constitutive law for wear, bearing friction or loss of actuator effectiveness.

The controller uses exact encoders and nominal computed torque:
\begin{equation}
 \tau_u=\bar M(q)\{\ddot q_d-K_D(\dot q-\dot q_d)-K_P(q-q_d)\}
       +\bar c(q,\dot q)+\bar g(q),\label{eq:controller}
\end{equation}
with $K_P=\omega^2I$, $K_D=2\omega I$. The standard nominal error equation and its sensitivity to model error are described in \cite{lynch}. Here actuation is ideal. A simulated torque transducer downstream of the actuator and upstream of the loading fixture reads $\tau_s=\tau_u+\xi_\tau$. A separate monitor receives the selected residual, task sign and time; the transducer noise is independent of feedback. Access to exact commands or complete robot data changes this information interface, as discussed in Section~\ref{sec:scope}.

Subtract nominal gravity at the measured configuration:
\begin{equation}
 r_\tau=\tau_s-\bar g(q)
       =w(q)\beta_m+e_fa+\xi_\tau+e_{\rm dyn},\quad
 w=\partial g/\partial m,\label{eq:static}
\end{equation}
where $e_{\rm dyn}=M\ddot q+c+\dot m A\dot q$. At rest $e_{\rm dyn}=0$; finite-time admission uses the bound below. The information supplied to this monitor defines the statistical experiment.

\subsection{Mirrored gravity and the selected observation}
For a planar chain with axial centers of mass, cumulative angles $\phi_j=\sum_{i\le j}q_i$ give a potential that is a weighted sum of $\sin\phi_j$. Reflection
\begin{equation}
 \mathcal Rq=(\pi-q_1,-q_2,\ldots,-q_n)
\end{equation}
leaves the potential unchanged and has $D\mathcal R=-I$. Thus $g(\mathcal Rq)=-g(q)$ and $w(\mathcal Rq)=-w(q)$. At the two reference poses, the faulted joint's channel $j$ has the scalar budgets
\begin{equation}
 b=|w_j(q_a)|b_m,\qquad r_d=|w_j(q_a)|r_m,
 \qquad\sigma^2=(\Sigma_\tau)_{jj}.
\end{equation}
Fixed errors in other reflection-compatible gravity parameters likewise become a common scalar offset after sign transformation. Their magnitude must be included in a finite $b$; using $B=\infty$ is an enlarged statistical class, not a certification of unbounded physical payload.

For complete-vector observations, whiten and split along $g=\Sigma_\tau^{-1/2}w(q_a)$, with $c_f=g^T\Sigma_\tau^{-1/2}e_f/\|g\|$ and $F=\|\Sigma_\tau^{-1/2}e_f\|^2$. The unconfounded contribution is $(F-c_f^2)\sum_{i\in I_\pi}a_i^2/2$. Proposition~\ref{prop:rankone} accounts for both this term and the total missing-slot coefficient $F$. Section~\ref{sec:scope} discusses the information restriction and the unconfounded energy in this example.

\subsection{Certified transitions and information cost}
The certificate covers $|\beta_m|\le.02$ kg, $|\dot\beta_m|\le.001$ kg/s, $0\le f\le.75$ N m and $|\dot f|\le.003$ N m/s. Static mismatch is absorbed into the moving equilibrium
\begin{equation}
 d=-M_0^{-1}(\beta_mg_p+e_ff),\qquad
 q_e=q_d+\omega^{-2}d(q_e,\beta_m,f).
\end{equation}
At $\omega=40$, Appendix~\ref{app:embedding} gives an equilibrium displacement below $1.22\times10^{-3}$ rad, and, after a fixed 1.5 s hold,
\begin{align}
 \|\dot q\|&\le1.91\times10^{-5}\ \mathrm{rad/s},\nonumber\\
 \|\ddot q\|&\le6.12\times10^{-4}\ \mathrm{rad/s^2},\nonumber\\
 \|e_{\rm red}\|&\le2.031\times10^{-3}\ \mathrm{N\,m}.
 \label{eq:extended_hold}
\end{align}
The same bounds close the transition state tube and command limits. With a 1.5 s move and $\delta=.25$ s, the reference calendar uses $k=12$.

Appendix~\ref{app:embedding} gives, with $x=e^{-12t}$,
\begin{equation}
 E(t)\le E_\infty+C_1x+C_2x^2\le E_\infty+Cx,
 \quad C=C_1+C_2.\label{eq:admission_constants}
\end{equation}
The exact rational coefficients have values $E_\infty\approx.002026081608$, $C_1\approx310.705171$, $C_2=7.43563476$, and $C\approx318.140806$, in N m.
For $\varepsilon>E_\infty$, the simple envelope gives
\begin{equation}
 k_{\min}(\varepsilon)=\left\lceil\frac{T_{\rm move}
 +\lambda_s^{-1}\pos{\log[C/(\varepsilon-E_\infty)]}}{\delta}\right\rceil,
 \quad\lambda_s=12.\label{eq:k_control}
\end{equation}
At $\varepsilon=.002031$ N m the hold bound is $1.49875$ s, so $k_{\min}=12$. At $\varepsilon=.01$ N m it is $.88284$ s and $k_{\min}=10$. The experiments in Sections~\ref{sec:slow}--\ref{sec:full_disclosure} retain $k=12$ and the tighter polynomial-envelope radius $.002030814$ N m. The smaller target is only $0.243\%$ above $E_\infty$; halving that gap increases the hold requirement by $\log2/12\approx.0578$ s and the integer budget to 13. This sensitivity concerns the certificate, not observed mechanical fragility.

A larger admission tolerance may shorten a transition, but also enlarges the uncertainty in the statistical test. This trade-off must be evaluated after recomputing the calendar. For a fixed calendar its protected power decreases with the tolerance.
\begin{proposition}[Offline admission--calendar design]\label{prop:controller}
Fix a controller, a known horizon $H$, a common noise variance, and a finite set of feasible exclusion budgets $\mathcal I$. For each $k\in\mathcal I$, set $t_k=k\delta-T_{\rm move}\ge0$ and fix a calendar $\pi_{H,k}$ before observing data. Assume a common admitted prefix and a verified envelope $E_\infty+Ce^{-\lambda_st}$. Define
\begin{equation}
 e_k=E_\infty+Ce^{-\lambda_st_k},\quad
 v_k=s_k-\operatorname{proj}_{\K_{\pi_{H,k}}}s_k,
\end{equation}
and restrict tolerances to $e_k\le\varepsilon\le\varepsilon_{\max}$. For $v_k\ne0$ put
\begin{equation}
 \mathcal P_k(\varepsilon)=\max\!\left\{\alpha,
 \Phi\!\left(\frac{\|v_k\|}{\sigma}
 -\frac{2\varepsilon\|v_k\|_1}{\sigma\|v_k\|}
 -z_{1-\alpha}\right)\right\};\label{eq:admission_objective}
\end{equation}
for $v_k=0$ set $\mathcal P_k=\alpha$. Then the best certified power within this calendar and fixed-direction test family is
\begin{equation}
 \max_{\substack{k\in\mathcal I,\ e_k\le\varepsilon\le\varepsilon_{\max}}}
 \mathcal P_k(\varepsilon)
 =\max_{\substack{k\in\mathcal I\,:\ e_k\le\varepsilon_{\max}}}
 \mathcal P_k(e_k).\label{eq:admission_opt}
\end{equation}
The maximum with $\alpha$ allows a constant randomized test as a fallback. All choices use model quantities, not observations.
\end{proposition}
\begin{proof}
Under the stated envelope-based admission rule, $(k,\varepsilon)$ is feasible when its declared tolerance satisfies $\varepsilon\ge e_k$. At a nonzero fixed direction, the nonconstant term in \eqref{eq:admission_objective} has derivative $-2\|v_k\|_1\varphi(d_k)/(\sigma\|v_k\|)<0$, where $d_k$ is its displayed normal argument and  $\varphi$ is the standard normal density. Taking the maximum with a constant preserves nonincrease. Hence the smallest feasible declared tolerance $e_k$ is optimal for each $k$; the remaining finite maximum is attained. Each selected test is level $\alpha$ because its error support is included on both sides. Deterministic selection before the experiment preserves that guarantee.\end{proof}
The result optimizes a lower guarantee within the stated family, not the enlarged-model minimax power. Fixed-calendar power is nonincreasing in $\varepsilon$; changing $k$ creates the discrete trade-off through both observations and $v_k$. For any feasible $k$, Corollary~\ref{cor:delay} gives the ideal time correction $(2\sqrt2/5)\sqrt{k\delta R/\kappa}\sqrt{T_\Lambda}$. Different controller gains require new stability, error and effort certificates.

\subsection{Error-protected statistical implementation}
For $v=s-z^*$ and a proved error bound $\varepsilon$ under each hypothesis, a valid threshold adds $\varepsilon\|v\|_1$ to \eqref{eq:implement_test}. Its guaranteed power is
\begin{equation}
 \underline{\mathcal B}_\pi=
 \Phi\!\left(\frac{\|v\|^2-2\varepsilon\|v\|_1}{\sigma\|v\|}
 -z_{1-\alpha}\right),\qquad v\ne0.
 \label{eq:guard_power}
\end{equation}
The two error sets are distinct and each pays its budget. This is the worst-power guarantee for the fixed ideal separating direction under an outer error box, not the exact minimax power of the expanded error model. When $G=0$, the two ideal mean classes intersect, so their minimax power is $\alpha$, attained by a constant randomized test; \eqref{eq:guard_power} is not evaluated at $v=0$. Nonconstant level-$\alpha$ tests may also exist. The numerical section compares this fallback and the nonzero-score guarantees. Readout variances and the $.25$ s independent sampling interval are specified model inputs; the physical calibration requirements are collected in Section~\ref{sec:scope}.

\section{Fixed-History Mechanisms}
\begin{proposition}[Three explicit projections]\label{prop:history}
For a fixed calendar, $G=0$ exactly when $s\in\K_\pi(B,r)$. On post-onset slots $1,\ldots,d$, with no prefix, put $v_0=\min(\eta,r)$. If $B\ge\eta(d+1)/2+v_0(d-1)/2$, then $s_j=\eta j$ has projection
\begin{equation}
 z_j^*=\eta\bar j+\min(\eta,r)(j-\bar j),\quad\bar j=(d+1)/2,
\end{equation}
and information
\begin{equation}
 G_d=\frac{(\eta-r)_+^2d(d^2-1)}{24\sigma^2}.\label{eq:linear}
\end{equation}
With constant nuisance ($r=0$), $B\ge\rho d/(n_0+d)$, $n_0$ zero-target prefix observations and $d$ subsequent observations of height $\rho$,
\begin{equation}
 G=\frac{\rho^2n_0d}{2\sigma^2(n_0+d)}.\label{eq:prefix}
\end{equation}
For two observations of target $(0,\rho)$ separated by $g_s$ slots, $B\ge\rho$,
\begin{equation}
 G=\frac{(\rho-rg_s)_+^2}{4\sigma^2}.\label{eq:gap}
\end{equation}
\end{proposition}
\begin{proof}
The zero-information assertion is membership in a closed convex set. If $\eta\le r$, the linear target is feasible. If $\eta>r$, for any feasible $z$ the sequence $z_j-rj$ is nonincreasing. Its inner product with the increasing, zero-sum sequence $j-\bar j$ is nonpositive; this is the projection variational inequality for the stated $z^*$. Summing $(j-\bar j)^2=d(d^2-1)/12$ proves \eqref{eq:linear}. For \eqref{eq:prefix}, subtract the mean of the $n_0+d$ entries. For \eqref{eq:gap}, the closest pair is centered at $\rho/2$ with separation $\min(\rho,rg_s)$.\end{proof}

For fixed finite $B$, a growing target eventually exceeds the box even if $\eta\le r$. The rate condition alone therefore does not decide whether moving is necessary. With a fixed prefix and $B=\infty$, \eqref{eq:linear} gives the leading no-movement information: its ratio to the full-calendar oracle tends to $(\eta-r)_+^2/(4\eta^2)$. A constant fault can be distinguishable from the prefix even when its post-onset-only history is not.

For fixed sample target values and covariance, lengthening a gap only expands the health class. If $q_i$ is the signed multiplier of its difference constraint in the objective $\frac12\norm{s-z}^2$, the derivative of $G$ with respect to that gap, when it exists, is $-r|q_i|/\sigma^2$. For progressive targets the total time derivative also includes the changing fault values; the gap derivative alone is not a progressive-fault scheduling gradient.

\section{Constant-Amplitude Scheduling}
\begin{theorem}[Best long-run information rate]\label{thm:constant}
Let $a_j=\rho>0$, $B\ge\rho$, $r>0$, $k\ge1$, and $n_0$ be fixed. Define
\begin{align}
 A_\rho&=\pos{\rho-r(k+1)/2},\\
 e_i(n)&=\pos{A_\rho-r\min(i,n-1-i)},\\
 Q(n)&=\sum_{i=0}^{n-1}e_i(n)^2,\\
 \mathcal J_*&=\max_{n\ge1}\frac{Q(n)}{2\sigma^2\delta(n+k)}.\label{eq:constant_rate}
\end{align}
Then
\begin{equation}
 \lim_{H\to\infty}\frac{\max_{\pi\in\PiH}G_H^B(\pi)}{(H+n_0)\delta}
 =\mathcal J_*.
\end{equation}
A repeated equal-hold calendar $\pi_*$ satisfies
\begin{equation}
 G_H^B(\pi_*)\ge\mathcal J_*(H+n_0)\delta-C
\end{equation}
with a constant independent of $H$. If $A_\rho>0$, an even maximizer exists in $1\le n\le2\lceil A_\rho/r\rceil$. The rate increases from $n$ to $n+1$ precisely when
\begin{equation}
 (n+k)[Q(n+1)-Q(n)]\ge Q(n).\label{eq:marginal}
\end{equation}
The sign of this comparison changes at most once.
\end{theorem}
\begin{proof}
The full cyclic KKT, all-calendar upper construction, and matching open-history dual are given in Appendix A. Those arguments retain the prefix and establish the bounded gap.\end{proof}
For fixed $k,\rho$ and $r\downarrow0$, \eqref{eq:marginal} yields $n_*\sim\sqrt{2k\rho/r}$. The information saturation length is instead of order $\rho/r$. For $\rho=.5,r=.1,k=2,\delta=1,\sigma=.4$, the optimal rate is $37/192$ at $n=4$; waiting until saturation at $n=8$ gives $21/160$, a $46.8\%$ rate difference. This example illustrates the marginal rule rather than establishing it. For $\mathcal J_*>0$, the corresponding long-horizon planning approximation is $H_{\rm plan}\approx\Lambda_{\alpha,\beta_{\rm miss}}/(\mathcal J_*\delta)$ slots, or $T_{\rm plan}\approx\Lambda_{\alpha,\beta_{\rm miss}}/\mathcal J_*$ seconds. Near the finite-prefix or zero-rate regime, use the exact projection instead.

For comparison, when $r=0$ and $B\ge\rho$, the exact finite-$H$ optimum is the better of staying and one balanced switch. Let $n=H-k>0$ and $p_n=n\bmod2$. The two informations are
\begin{equation}
 \frac{\rho^2Hn_0}{2\sigma^2(H+n_0)},\qquad
 \frac{\rho^2}{2\sigma^2}\left(n-\frac{p_n}{n_0+n}\right).\label{eq:R0}
\end{equation}
Indeed, subtracting the full-history mean leaves $\rho^2[n-(n_+-n_-)^2/(n_0+n)]$. Balance minimizes the second term, one switch can realize it, and further switching reduces $n$. When $n$ is even, moving is strictly preferable exactly when $k<H^2/(n_0+H)$; the parity term in \eqref{eq:R0} is required otherwise.

\section{Sharp Loss for a Progressive Fault}
We now take $a_j=\rho_0+\eta j$ with $\eta>0$. The full-calendar oracle is cubic, so first-order efficiency alone conceals the scheduling cost.

\subsection{A horizon-dependent deterministic construction}
Choose $\xi>0$. The achieving construction takes transitions exactly $k$ slots long, a feasible subclass of Definition~\ref{def:calendar}. A symmetric block has an even $n=2m$ and comprises $m$ readings at $+$, a $k$-slot transition, $n$ readings at $-$, a second transition, and $m$ readings at $+$. Its duration is $L=2(n+k)$ slots. Blocks concatenate at task $+$ without an additional transition.

For a planned $H$, first select as many blocks with $n_\ell=2\lceil\xi\ell/2\rceil$ as fit. Distribute the remaining multiples of four slots across these blocks by increasing their $n_\ell$ by even integers as evenly as possible. Fewer than four residual slots are observed at $+$ with zero weight in the constructive certificate. If no block fits, stay. Denote this fully specified calendar by $\pi_H^\xi$. It depends on the planned endpoint and parameters, never on data.

\begin{theorem}[Sharp second-order deficit]\label{thm:sharp}
Fix $\rho_0\ge0$, $\eta,r>0$, $k\ge1$, $\sigma>0$ and finite $n_0$. With $B=\infty$,
\begin{equation}
 \boxed{D_H^{\infty,*}
 =\frac{\sqrt2}{5\sigma^2}\sqrt{kr}\,\eta^{3/2}H^{5/2}
 +o(H^{5/2}).}\label{eq:sharp}
\end{equation}
The calendar $\pi_H^\xi$ attains this coefficient at $\xi_*=2k\eta/r$. For every fixed $\xi>0$ its deficit obeys
\begin{equation}
 2\sigma^2 D_H^{\infty}(\pi_H^\xi)
 \le\frac15\left(\frac{2k\eta^2}{\sqrt\xi}+r\eta\sqrt\xi\right)H^{5/2}
 +O(H^2).\label{eq:upper}
\end{equation}
\end{theorem}
\begin{proof}
Appendix B first proves the schedule-uniform converse coefficient $2\sqrt2/5$ in the normalization $2\sigma^2D_H/(\sqrt{kr}\eta^{3/2}H^{5/2})$. It then proves \eqref{eq:upper} with a local exact projection residual on each symmetric block. Their combination at $\xi=2k\eta/r$ gives \eqref{eq:sharp}.\end{proof}

The local mechanism is useful for implementation. With $c_0=(k+1)/2$ and center amplitude $a$, the local constant-target projection has levels $r(c_0+j)$ with the signs prescribed in Appendix B. Its loss is
\begin{equation}
 L_c(a,n)=ar\,n(n+2k)-4r^2\sum_{j=0}^{n/2-1}(c_0+j)^2.\label{eq:Lc}
\end{equation}
The projection residual $w$ has zero sum and satisfies $\sum_iw_i\epsilon_i(t_i-c)=0$. Thus it can be used for the affine target without changing its dual pairing. Concatenation refers to these \emph{dual directions}, not to the local primal optimizers. The only support relation needed is
\begin{equation}
 h_{\K_{\rm whole}}(w)\le\sum_\ell h_{\K_\ell}(w_\ell),\label{eq:support}
\end{equation}
which follows by restricting every global feasible path to a block. Adjacent local primal optimizers need not themselves join within the rate budget.

For blocks satisfying $a_\ell\ge r(k+n_\ell-1)/2$, there is a closed finite-$H$ information lower bound without solving a QP:
\begin{equation}
 G_H^\infty(\pi_H^\xi)\ge\frac1{2\sigma^2}
 \sum_{\ell\ \mathrm{admissible}}
 \bigl[2n_\ell a_\ell^2-L_c(a_\ell,n_\ell)\bigr].\label{eq:closedG}
\end{equation}
Zero-weight early blocks and the prefix remain in the actual experiment.

\begin{remark}[Tuning, necessity, and planned horizons]
At $\xi=\lambda\xi_*$, the proved upper coefficient relative to the sharp one is
\begin{equation}
 \mathcal R(\lambda)=\tfrac12(\lambda^{-1/2}+\lambda^{1/2}).\label{eq:robust}
\end{equation}
A factor-two schedule tuning error gives $\mathcal R\approx1.061$ and a factor-four error gives $1.25$. This bounds calendar tuning loss; it does not validate a detector using an incorrect fault template. For regular families $n(t)\asymp t^a$, a local balance suggests loss orders $H^{3-a}$ and $H^{2+a}$, and hence the exponent $a=1/2$. This is a design heuristic; the proved statements are the schedule-uniform converse and the achieving construction in Theorem~\ref{thm:sharp}, not a classification of every efficient run sequence.

The construction has $n(t)\sim\sqrt{2k a(t)/(R\delta)}$, the constant-amplitude rule evaluated at the local growing amplitude. The sharp planned-horizon coefficient is proved for terminal-balanced calendars. An anytime rule with an unfinished last block retains the known $O(H^{5/2})$ guarantee; its sharp coefficient at arbitrary stopping times is not asserted.
\end{remark}

\begin{corollary}[Second-order planned detection time]\label{cor:delay}
For fixed $\beta_{\rm miss}\in(0,1)$ define
\begin{equation}
 H_{\alpha,\beta_{\rm miss}}^*=
 \min\{H:\exists\pi\in\PiH,\ \mathcal B_\pi(\alpha)\ge1-\beta_{\rm miss}\}.
\end{equation}
Define the exact Gaussian requirement and its cubic reference horizon by
\begin{align}
 \Lambda_{\alpha,\beta_{\rm miss}}&=\tfrac12
 (z_{1-\alpha}+z_{1-\beta_{\rm miss}})^2,\nonumber\\
 H_\Lambda&=(6\sigma^2\Lambda_{\alpha,\beta_{\rm miss}}/\eta^2)^{1/3}.
\end{align}
Under Theorem~\ref{thm:sharp}, as $\alpha\downarrow0$,
\begin{equation}
 H_{\alpha,\beta_{\rm miss}}^*
 =H_\Lambda+\frac{2\sqrt2}{5}\sqrt{kr/\eta}\,H_\Lambda^{1/2}
 +o(H_\Lambda^{1/2}).\label{eq:delay}
\end{equation}
In physical time, $T_\Lambda=\delta H_\Lambda=[6\sigma^2\delta\Lambda_{\alpha,\beta_{\rm miss}}/\kappa^2]^{1/3}$ and the correction is
$(2\sqrt2/5)\sqrt{\tau_mR/\kappa}\sqrt{T_\Lambda}+o(\sqrt{T_\Lambda})$.
\end{corollary}
\begin{proof}
Equation~\eqref{eq:power} requires $G_H^*\ge\Lambda_{\alpha,\beta_{\rm miss}}$. Substitute $G_H^*=\eta^2H^3/(6\sigma^2)-C H^{5/2}+o(H^{5/2})$ and set $H=H_\Lambda+u\sqrt{H_\Lambda}$. Matching the $H_\Lambda^{5/2}$ terms gives $u=C/[3\eta^2/(6\sigma^2)]$. The oracle's quadratic term and rounding are lower order. Optimal information is nondecreasing in $H$, which justifies inversion. The exact threshold improves finite-tail use without turning the asymptotic expansion into a finite-horizon error bound.\end{proof}
The penalty is a square-root function of healthy drift during one transition relative to fault growth over the observation horizon. This is a high-probability planned-time statement, not a mean alarm delay or ARL result. Summable testing levels, for example proportional to $\alpha/(H+1)^2$, control anytime false alarms for a single predetermined calendar and preserve its first-order time law. Such a scan does not physically implement all horizon-dependent calendars simultaneously.
For unknown onset, a fixed-hold periodic calendar is an onset-independent default. For each fixed onset and finite hold $n$, a positive affine growth eventually retains the observed-slot fraction $n/(n+k)$ of the full-calendar oracle in this ideal path model; the finite initial phase changes only boundary terms. A scan over onset and length uses one physically fixed calendar and summable levels. Neither this observation nor alpha spending transfers the sharp planned-horizon coefficient to an adaptive or unknown-onset design.

\section{Fixed Finite Healthy Bounds}
\begin{theorem}[Bounded switching advantage]\label{thm:finite}
Fix $B,r,\eta>0$, $k\ge1$, $\rho_0\ge0$ and finite $n_0$. Define the full polynomial
\begin{align}
 P_B(H)&=\frac1{2\sigma^2}\sum_{j=1}^H(2Ba_j-B^2)\nonumber\\
 &=\frac{B\eta}{2\sigma^2}H^2+
 \frac{2B\rho_0+B\eta-B^2}{2\sigma^2}H.\label{eq:PB}
\end{align}
Then
\begin{equation}
 D_H^{B,*}=P_B(H)+O(1),\label{eq:finite}
\end{equation}
and a parameter-dependent constant $C_B$, independent of $H$, satisfies
\begin{equation}
 0\le\max_{\pi\in\PiH}G_H^B(\pi)-G_H^B(\mathrm{stay})\le C_B.\label{eq:regret}
\end{equation}
\end{theorem}
\begin{proof}
Appendix C constructs one capped rate-limited healthy tracker over the full history, preserving the free terminal value. Its possible losses relative to $P_B$ have a finite total independent of the calendar. The no-movement projection has the same polynomial up to a bounded early correction.\end{proof}

For any fixed calendar, $\K_\pi(B,r)\subseteq\K_\pi(\infty,r)$ implies $G_H^B(\pi)\ge G_H^\infty(\pi)$ and $D_H^B(\pi)\le D_H^\infty(\pi)$. Hence the constructive finite-$H$ bound from the unbounded-level experiment is conservative for the same finite-box calendar; the asymptotic sharp coefficient alone is not a finite-$H$ upper error bar. Its inverse bound is specific to the unbounded class; the finite-box inverse is Theorem~\ref{thm:finite}.

The result allows a finite, potentially large benefit from early switching. It neither identifies the earliest optimal stopping amplitude nor makes staying optimal for every finite experiment. An informative local comparison is $\sqrt{2kr}\,a^{3/2}$ versus $2Ba$, which suggests the amplitude scale $a\asymp B^2/(kr)$. That comparison is a design heuristic, not an exact crossing threshold. Integrating the local saving balance suggests the dimensional proxy $B^5/(k^2r^2\eta\sigma^2)$ for transient advantage. It is neither the uniform constant $C_B$ nor a certified finite-horizon bound. In particular, the finite-$B$ and $B\to\infty$ asymptotics cannot be interchanged. With $r=0$, a constant healthy parameter remains shared indefinitely, so Theorem~\ref{thm:finite} is inapplicable.

\begin{table}[t]
\caption{Design conclusions within the specified experiment}
\label{tab:design}\centering\small
\begin{tabular}{p{.22\columnwidth}p{.68\columnwidth}}
\toprule
Regime & Consequence\\\midrule
Constant fault, $r=0$ & Stay or one balanced switch, with the exact parity correction in \eqref{eq:R0}.\\
Constant fault, $r>0$ & Repeat a maximizer of \eqref{eq:constant_rate}; switch by marginal information, not saturation.\\
Growing fault, $B=\infty$ & Symmetric growing-hold blocks attain the sharp loss and \eqref{eq:delay}.\\
Growing fault, fixed $B$ & Switching may be valuable over the intended finite horizon; eventually staying has the exact quadratic baseline. Choose a stopping stage by finite-calendar information, not a universal threshold.\\\bottomrule
\end{tabular}
\end{table}


\section{Numerical Illustration}
We compare schedules by whole-history projection and protected power. The mechanical setting uses $B=.03313$ N m, $r=4.141\times10^{-4}$ N m/slot, $k=12$, $\delta=.25$ s and four prefix readings. The primary slow-ramp example uses the same specified readout level $\sigma=.01$ N m as the closed-loop illustration; additional noise levels display sensitivity. Physical interpretation and calibration conditions are stated once in Section~\ref{sec:scope}.

\subsection{Slow growth: switching before amplitude escape}\label{sec:slow}
The bias slope is $\kappa=3\times10^{-4}$ N m/s, giving $\eta=7.5\times10^{-5}<r$. The stay target belongs to the finite healthy set through $H=441$, so its minimax power is only $\alpha$ throughout that interval. In the unbounded-level class its information remains zero for every $H$. The symmetric schedule has $\xi=2k\eta/r\approx4.35$ and a first complete block of 36 slots. At $H=300$ (75 s), it uses ten moves and has protected power above $0.999$ at $\alpha=.001$ and $\sigma=.01$, while stay has zero information. This demonstrates useful task contrast under the analytic error radius, before amplitude escape at about 110 s.

\begin{table}[t]\centering\small
\caption{Slow-ramp powers, $\alpha=.001$. $P_i$ is ideal power; $P_g$ uses the same per-side error radius $.002030814$ N m.}\label{tab:slow}
\setlength{\tabcolsep}{2.5pt}
\begin{tabular}{rrrrrrr}\toprule
$H$&$\sigma$&Moves&Stay $P_g$&Fixed-76 $P_g$&Sym. $P_i$&Sym. $P_g$\\\midrule
200&0.01&8&$\alpha^\dagger$&0.065&0.997&0.301\\
300&0.01&10&$\alpha^\dagger$&0.997&$>0.999$&$>0.999$\\
300&0.05&10&$\alpha^\dagger$&0.027&0.301&0.080\\
500&0.05&14&0.0006&0.851&$>0.999$&0.985\\
600&0.10&16&0.004&0.492&0.914&0.749\\
800&0.10&20&0.208&0.997&$>0.999$&$>0.999$\\
1000&0.20&24&0.259&0.941&0.985&0.956\\
\bottomrule\end{tabular}
\par\vspace{2pt}\parbox{.98\columnwidth}{\footnotesize Moves refer to the symmetric schedule. $\dagger$: $G=0$; a constant randomized test rejects with probability $\alpha$. The nonzero-score protection formula is not evaluated. Other rows report the particular protected score, even when its bound falls below $\alpha$.}
\end{table}
The fixed-76 rule is a simple comparator obtained from the local constant-amplitude scale at $a=.1$ N m. It is not refitted at each horizon. At $H=300,\sigma=.01$, it also succeeds, with protected power $.9967$: task switching, rather than the exact growing-hold law alone, explains the main zero-information contrast. The growing-hold rule buys more at intermediate horizons. At $H=500,\sigma=.05$, its lower power is $.9846$ versus $.8514$ for fixed-76. At $H=800,\sigma=.1$, the respective values are $.9996$ and $.9974$; at $H=1000,\sigma=.2$, they are $.9564$ and $.9413$. Thus the protected-power advantage is concentrated in the intermediate tested range; saturation of power does not imply equal asymptotic information efficiency.

Early failures remain relevant. At $H=200,\sigma=.01$ the symmetric ideal power is $.9968$ but protection reduces it to $.3011$. At $H=300,\sigma=.05$ its protected power is only $.0797$. Neither the small pointwise residual nor switching by itself guarantees the requested power. The complete calculation covers eight horizons, three policies, two healthy classes, four noise levels and two test levels. These horizons reach at most 250 s and $.075$ N m, within the analytic envelope of Appendix~\ref{app:embedding}.
\begin{figure}[t]\centering
\includegraphics[width=\columnwidth]{figures/08_slow_power_comparison.pdf}
\caption{Slow-ramp powers at $\sigma=.05$ and $\alpha=.001$. The fixed-76 baseline and the symmetric protected score use identical model-error and health budgets. The ideal symmetric curve separates statistical signal from implementation protection.}\label{fig:slowpower}
\end{figure}

\subsection{Selecting an admission tolerance}
Proposition~\ref{prop:controller} is evaluated at fixed gain 40 over exclusion budgets $k=7,\ldots,24$, with tolerance cap $.01$ N m. The envelope rules out $k=7,8,9$; all budgets $10,\ldots,24$ are feasible. For each budget the symmetric calendar and its whole-history projection are recomputed before data collection. A common 4.5 s initial hold precedes the prefix, so all candidates use an already admitted prefix. This common setup is distinct from the post-onset budget $H$.
\begin{table}[t]\centering\small
\caption{Admission design at $H=500$, $\sigma=.05$, $\alpha=.001$. The radius $e_k$ is the simple exponential bound.}\label{tab:admission}
\setlength{\tabcolsep}{4pt}
\begin{tabular}{rrrrr}\toprule
$k$&Hold (s)&$10^3e_k$ (N m)&Moves&$\mathcal P_k(e_k)$\\\midrule
10&1.00&3.9808&16&0.8843\\
11&1.25&2.1234&16&0.9874\\
12&1.50&2.0309&14&0.9846\\
14&2.00&2.0261&14&0.9736\\
\bottomrule\end{tabular}
\end{table}
The best of the fifteen feasible budgets is $k=11$: the shorter wait compensates for its larger error radius. Increasing $k$ beyond 12 barely changes the floor-dominated error but consumes more observation time. At $H=300,\sigma=.05$, the same best budget raises the bound only from $.0797$ to $.0862$, still far below $.9$. A joint admission choice is a useful finite design operation, not a remedy for an inadequate signal. The baseline results in Table~\ref{tab:slow} remain at $k=12$; the new optimum does not replace them.
\begin{figure}[t]\centering
\includegraphics[width=\columnwidth]{figures/09_admission_tradeoff.pdf}
\caption{The complete feasible admission comparison for $H=500$, $\sigma=.05$, $\alpha=.001$. The optimum is over the stated finite calendar family, not all feedback laws, tolerances or schedules.}\label{fig:admission}
\end{figure}

\subsection{Fast growth: a small residual norm is insufficient}\label{sec:full_disclosure}
A separate faster-ramp example uses $\kappa=.003$ N m/s with the same health and implementation budgets. Table~\ref{tab:guard} includes the low-noise failures and representative higher-noise endpoints. All twelve original ideal first-success records fall below $.9$ after protection; the complete records are given in Supplementary Table S1. Coincident stay/symmetric calendars are grouped here. The low-noise outcomes expose the effect of the coherent penalty $2\varepsilon\|v\|_1/(\sigma\|v\|)$, not merely the ratio $\varepsilon/\sigma$.
\begin{table}[t]\centering\small
\caption{Fast-ramp ideal endpoints: implementation protection. All twelve endpoint records fail the protected $.9$ target.}\label{tab:guard}
\setlength{\tabcolsep}{4pt}
\begin{tabular}{lrrrrr}\toprule
Policy&$\sigma$&$H$&$\alpha$&$P_i$&$P_g$\\\midrule
stay/sym.&0.01&56&$10^{-3}$&0.908&0.079\\
stay/sym.&0.01&70&$10^{-8}$&0.930&0.082\\
sym.&0.10&119&$10^{-3}$&0.905&0.835\\
sym.&0.20&183&$10^{-3}$&0.904&0.861\\
\bottomrule\end{tabular}
\end{table}
These lower guarantees concern a specified linear score under an outer error box; they need not equal its rejection probability on a particular mechanical path. The successful slower-ramp cases and all of these failures use the same analytic error radius, rather than a trajectory-dependent reduced radius.

\subsection{Sharp limits and bounded transient advantage}
For $r=\eta=.1$, $k=2$, $n_0=4$ and $B=\infty$, the exact normalized symmetric-calendar deficits at $H=200,400,800,1600$ are $.5601,.5610,.5623,.5630$, approaching $2\sqrt2/5\approx.5657$. The finite certificate recovers $98.96\%,99.48\%,99.75\%,99.87\%$ of each calendar's exact information. These are specified-calendar projections, not finite-horizon global optima. The constant-amplitude example gives rates $37/192$ and $21/160$ for the marginally optimal hold and the saturation hold. The exact-threshold planned-time approximation gives $124.6,141.1,164.5,189.4$ slots versus exact first successful endpoints $126,142,165,190$ in the separate mathematical example. Gaussian power and preassigned multiple-test levels have finite Monte Carlo implementation checks; the analytical formulas establish the guarantees.
\begin{figure}[t]\centering
\includegraphics[width=\columnwidth]{figures/01_sharp_convergence.pdf}
\caption{Exact information for the symmetric calendar, its finite constructive bound, and the sharp limiting coefficient. The theorem concerns the optimum over deterministic calendars.}\label{fig:sharp}
\end{figure}

For the fast-ramp robot scaling at $\sigma=.2$, the uniform bound in \eqref{eq:explicit_CB} is about $2.97\times10^5$ nats. The dimensional proxy $B^5/(k^2r^2\eta\sigma^2)$ is $53.9$ nats, and the best tested switch-then-stay candidate gains $14.7$ nats over stay. These three values quantify, respectively, a proved calendar-uniform bound, an unproved local-balance scale, and a finite candidate comparison. The first is loose because the proof charges every boundary up to $N_B$ worst-case deficiencies of size $4Ba_j$. The finite-box theorem limits growth of the benefit, not its practical magnitude.

\subsection{Closed-loop reduction checks}
The six-coordinate rigid mechanism was tested with 16 event-aligned extreme paths, 240 random bounded-rate paths and six long histories, in addition to the short illustrative traces. The long histories reach $.75$ N m over 250 s. Their largest admitted residual-vector norm is $1.34\times10^{-4}$ N m, versus the analytic $2.03\times10^{-3}$ N m bound. The latter, not the observed maximum, is used in all protected scores. Noisy-readout calculations and the finite nonlinear stress tests serve different purposes: the former evaluate the stated statistical experiment and the latter check its implementation.

For a 19-reading calendar, the ideal selected-channel information is $32.76$ and the minimax power at $\alpha=10^{-8}$ is $.9935$. A specific admissible path has $.9994$ rejection probability; complete-vector information is $44.27$. Those values demonstrate both the distinction between pathwise and worst-class performance and the additional information available to a complete-vector monitor.

\section{Scope and Implementation Decisions}\label{sec:scope}
The result is a planned assessment of a specified additive alternative at a declared onset and covariance. A torque-loading fixture can generate that alternative for a synchronized self-test.  A fixed periodic calendar can serve an unknown-onset scan with preassigned error allocation, but the present sharp planned-horizon theorem does not supply an ARL or Lorden optimum.

\paragraph{Observation information.}
The scalar experiment applies to a monitor supplied one torque residual, rather than to an estimator with all control signals. In the ideal-actuation simulator, exact command access gives a noiseless effort signal; all six effort channels also reveal an unconfounded fraction $1-\chi=94.24\%$ in the present geometry. Those interfaces should use their additional information. A full-rank nuisance matrix could remove the latter component, but its vector path class requires its own analysis. 

\paragraph{Unmodeled parity.}
The task contrast separates task-odd healthy effects from a task-even additive bias. Multiplicative actuator faults can reverse with gravity and remain confounded with load. Output-side static friction, sensor zero drift, and off-axis center-of-mass errors can contain task-even healthy components. Their amplitude and rate sets must be included before deploying a test; under the current scalar class they are excluded. For the planar map, a 2 cm off-axis offset produces a 30\% parity defect, not an admissible small-error validation. Mirror-compatible fixed offsets consume a finite amplitude budget or must be calibrated. Underestimating any health envelope can invalidate false-alarm protection.

\paragraph{Randomized and adaptive schedules.}
The converse chooses a health path after a deterministic calendar is specified. If health instead commits before a random calendar, the order changes: for path loss $L$,
\begin{equation}
 \inf_z\E_\pi L(\pi,z)\ \ge\ \E_\pi\inf_z L(\pi,z).
\end{equation}
For example, $(1-z)^2$ and $(-1-z)^2$, chosen with equal probability over $z\in[-1,1]$, give $1$ on the left and $0$ on the right. Thus neither the sharp constant nor its exponent is established here for that randomized game or for measurement-adaptive control. Classical active testing \cite{chernoff,naghshvar} provides the appropriate wider perspective.

\paragraph{Practical use.}
Specify the monitored channel and verify the parity of its healthy and fault regressors. Certify transition, effort and residual-error bounds and choose a sampling covariance. Use upper health bounds and a specified fault alternative to generate a deterministic constant-hold, symmetric or switch-then-stay calendar. Compute its full-history projection and guarded power, including any prefix and all transition exclusions. Accept a plan only if the guarded power meets the requirement; tune or lengthen the plan using this finite calculation rather than an asymptotic time alone. For a pure growth ray, a detector built at $\kappa_{\min}$ has nondecreasing guaranteed power as growth increases, but its calendar need not remain optimal. Choose an admitted error budget jointly with the exclusion duration using Proposition~\ref{prop:controller}; then evaluate the complete plan. The declared noise variances are synthetic readout scenarios. Hardware use requires both residual-error and noise-bandwidth calibration, including a justification for the independent sampling interval. All compared admission budgets require a common prefix already admitted at their error level.

\section{Conclusion}
Deterministic two-task self-tests exhibit two different optimal design regimes under one rate-limited healthy class. Unbounded levels incur an unavoidable deficit with a sharp five-halves-power coefficient and a matching symmetric construction. Fixed amplitude bounds reduce the deficit to a full quadratic baseline plus a bounded remainder. The control interface turns a verified settling envelope into an offline choice between time exclusions and residual-error protection. In the slow-ramp example, switching separates a fault that staying cannot distinguish before amplitude escape. Its largest protected-power improvement over the tested fixed hold occurs at intermediate horizons; the powers become close at the longer tested horizons, without making the two information-efficiency laws equivalent. Randomized and adaptive calendars with precommitted health paths, unknown onset, vector or mixed-parity uncertainty, and physical calibration of residual error and noise remain open directions.
\appendices
\section{Proof of the Constant-Amplitude Scheduling Theorem}
\subsection{Cyclic projection}
\begin{proof}[Periodic projection and marginal rule]
Consider $2n$ cyclic observations with target $y=(\rho\ind_n,-\rho\ind_n)$. Consecutive edge budgets are $r$ within a plateau and $r(k+1)$ at either task change. Write $A=A_\rho$ and $e_i=e_i(n)$. The proposed positive nodes are $z_i=\rho-e_i$ and the negative nodes their negatives. They lie in $[-\rho,\rho]\subset[-B,B]$; their internal slopes have magnitude at most $r$, and their boundary jump is $2(\rho-A)\le r(k+1)$.

Let $D$ be the cyclic forward-difference matrix and $c_n=\frac12\sum_{i=0}^{n-1}e_i$. Signed edge multipliers are
\begin{equation}
 \lambda_i=c_n+\sum_{j=0}^i(z_j-y_j),\quad i=0,\ldots,2n-1.
\end{equation}
The total residual is zero, so the final multiplier is $c_n$. Hence $z-y+D^T\lambda=0$. On the positive plateau the multipliers decrease through zero at its center; on the negative plateau they increase through zero. Every nonzero multiplier has the sign of the corresponding saturated slope. A flat center has zero multipliers on its unconstrained edges. Thus
\begin{equation}
 \lambda_i(Dz)_i=|\lambda_i|h_i
\end{equation}
for every edge budget $h_i$. The convex quadratic KKT conditions prove optimality and uniqueness of the nodes. The half-squared residual is $Q(n)$.

Put $Q(0)=0$. Its increments are $A^2,A^2,(A-r)_+^2,(A-r)_+^2,\ldots$, so $Q$ is concave as a discrete sequence and saturates at $n_{\rm sat}=2\lceil A/r\rceil$ when $A>0$. Comparing adjacent ratios gives \eqref{eq:marginal}. With $\Delta_n=Q(n+1)-Q(n)$, the comparison numerator $B_n=(n+k)\Delta_n-Q(n)$ obeys
\begin{equation}
 B_{n+1}-B_n=(n+k+1)(\Delta_{n+1}-\Delta_n)\le0.
\end{equation}
After saturation the numerator is fixed and the denominator increases. An odd-index ratio between two even ones is their denominator-weighted average, while $Q(2)/(k+2)>Q(1)/(k+1)$ for $k,A>0$. An even maximizer therefore exists in the stated finite interval.
\end{proof}

\subsection{An upper bound for every open history}
\begin{proof}[Global rate and finite information gap]
Define $q_* =\max_n Q(n)/(n+k)$ and
\begin{equation}
 C_0=(\lceil\rho/r\rceil+1)\rho^2.
\end{equation}
Decompose any post-onset observation history into runs of lengths $n_1,\ldots,n_M$. On each run place the signed version of $\rho-e_i(n_j)$. Adjacent boundary values have magnitude $\rho-A$ and opposite signs differ by at most $r(k+1)$, so the actual excluded intervals allow a continuous rate-limited connection. Additional movement intervals cannot invalidate reachability. The candidate has squared residual $\sum_jQ(n_j)$.

To approximate the healthy-prefix target using a feasible adversary, choose zero on the prefix and clip this candidate after onset to $[-rt,rt]$ in slot time. This is a choice for an upper-bound construction, not a healthy initial-state assumption. Maxima and minima of scalar $r$-Lipschitz functions are $r$-Lipschitz. The clipping affects at most $\lceil\rho/r\rceil+1$ observations and adds at most $C_0$ squared residual. Since at least $k(M-1)$ transition slots occur,
\begin{equation}
 2\sigma^2G_H^B(\pi)
 \le\sum_jQ(n_j)+C_0
 \le q_*(H+k)+C_0.\label{eq:allconstupper}
\end{equation}
This inequality evaluates one feasible whole-history path; it does not interchange a minimum and a sum.

Now repeat an optimal plateau length $n_*$. Let $Q_*=Q(n_*)$, $c_*=\frac12\sum_ie_i(n_*)$. A repeated $m$-cycle closed problem has half-squared loss $mQ_*$. The cyclic rate-constrained quadratic dual is
\begin{equation}
 \lambda^TDy-\tfrac12\norm{D^T\lambda}^2-\sum_ih_i|\lambda_i|.
\end{equation}
Remove its one final closing edge to obtain a feasible dual on the open observed history. The removed multiplier is $c_*$, the target jump is $2\rho$, and the budget is $r(k+1)=2(\rho-A)$ when $A>0$. Removing the incidence vector changes the squared norm by $-4Ac_*+2c_*^2$. The net dual change is consequently
\begin{equation}
 -2\rho c_*+2Ac_*-c_*^2+r(k+1)c_*=-c_*^2.
\end{equation}
It follows that the open post-onset objective is at least $mQ_*-c_*^2$. Adding the complete prefix can only increase this lower bound. At $H_m=2m(n_*+k)-k$,
\begin{equation}
 G_{H_m}^B(\pi_*)\ge(mQ_*-c_*^2)/\sigma^2.\label{eq:constlower}
\end{equation}
For arbitrary $H$, use $m=\lfloor(H+k)/(2(n_*+k))\rfloor$ complete pairs and observe at the current task for the remainder. Combining \eqref{eq:allconstupper} and \eqref{eq:constlower} yields
\begin{equation}
0\le G_H^{B,*}-G_H^B(\pi_*)\le
\frac{2Q_*+C_0+2c_*^2}{2\sigma^2},
\end{equation}
where $G_H^{B,*}=\max_\pi G_H^B(\pi)$. Division by $(H+n_0)\delta$ proves Theorem~\ref{thm:constant}. If $A=0$, the same upper construction has $q_*=0$ and nonnegative information supplies the zero-rate lower bound.
\end{proof}

\begin{proof}[Small-rate holding scale]
Before saturation, for $n=2m$,
\begin{equation}
 Q(n)=2\left[mA^2-Ar m(m-1)
 +\frac{r^2m(m-1)(2m-1)}6\right].
\end{equation}
For $n$ of order $r^{-1/2}$, the adjacent-rate comparison equals
$kA^2-Arn^2/2+o(1)$, uniformly on compact positive multiples of $r^{-1/2}$. Since $A\to\rho$, unimodality brackets all optimizers between $(1\pm\epsilon)\sqrt{2k\rho/r}$ for small $r$.\end{proof}

\section{Proof of the Sharp Progressive-Fault Limit}
\subsection{A schedule-uniform converse}
\begin{proof}[Converse in Theorem~\ref{thm:sharp}]
For any feasible healthy path chosen to be zero on the prefix,
\begin{equation}
 2\sigma^2D_H^\infty(\pi)\ge
 \sum_{j\notin I_\pi}a_j^2+
 \sum_{j\in I_\pi}(2\epsilon_j a_jz_j-z_j^2).\label{eq:adv}
\end{equation}
The inequality follows by evaluating one feasible point in the projection objective.

Consider $L$ consecutive post-onset slots with minimum amplitude $A>0$. Let $d$ slots be excluded and let the clipped observed-run lengths be $n_1,\ldots,n_m$. Then $m\le d/k+1$. On every run use the zero-endpoint signed tent $z_i=\epsilon r\min(i,n-1-i)$. All pieces join to zero across block boundaries and excluded intervals. Set
\begin{equation}
 C=\frac r4\left(2A-\frac{rL}{2}\right).
\end{equation}
When $C\ge0$, tent height is at most $rL/2$, and
$\sum_{i=0}^{n-1}\min(i,n-1-i)\ge(n^2-2n)/4$. The local contribution to \eqref{eq:adv} is at least
\begin{align}
 A^2d+C\sum_i n_i^2-2C(L-d)
 &\ge A^2d+\frac{Ck(L-d)^2}{d+k}-2CL.\label{eq:macro1}
\end{align}
The right side remains valid when there are no observations. Put $y=d+k$ and apply AM--GM to
\begin{equation}
 (A^2+Ck)y+\frac{Ck(L+k)^2}{y}
 -2Ck(L+k)-A^2k-2CL.
\end{equation}
Thus every calendar has a feasible local explanation giving at least
\begin{equation}
 \begin{split}
 \Psi(A,L)={}&2(L+k)\sqrt{Ck(A^2+Ck)}\\
 &-2Ck(L+k)-A^2k-2CL.
\end{split}\label{eq:macro}
\end{equation}
The nonnegative part may be used, or zero when $C<0$.

Fix $\epsilon\in(0,1)$ and partition the late interval $[\lceil\epsilon H\rceil,H]$ into blocks of length $L=\lfloor H^{3/4}\rfloor$. Uniformly there, $A=\Theta(H)$ and $rL/A=o(1)$, so $C=(rA/2)(1+o(1))$. The leading sum in \eqref{eq:macro} is
\begin{equation}
 \sqrt{2kr}\sum_{\rm blocks}A^{3/2}L.
\end{equation}
The accumulated $A^2k$ terms are $O(H^3/L)=O(H^{9/4})$, the $CL$ and $CkL$ terms are $O(H^2)$, and replacing $C$ by $rA/2$ incurs $O(H^{9/4})$. The last incomplete block contributes only a lower-order remainder. The Riemann sum is
\begin{equation}
 \sum A^{3/2}L
 =\tfrac25\eta^{3/2}(1-\epsilon^{5/2})H^{5/2}+o(H^{5/2}).
\end{equation}
These bounds are uniform in the calendar because only its run-count inequality was used. Taking the infimum deficit before the limit and then letting $\epsilon\downarrow0$ gives
\begin{equation}
 \liminf_{H\to\infty}
 \frac{2\sigma^2D_H^{\infty,*}}{\sqrt{kr}\eta^{3/2}H^{5/2}}
 \ge\frac{2\sqrt2}{5}.\label{eq:sharp_lower}
\end{equation}
The converse remainder used here is $o(H^{5/2})$.
\end{proof}

\subsection{Symmetric local projection and affine cancellation}
\begin{proof}[Constructive upper bound]
For an even $n=2m$, let $L=2(n+k)$ and $c=(L+1)/2$. Local admitted times are
\begin{equation}
\begin{split}
 I={}&\{1,\ldots,m\}\cup\{m+k+1,\ldots,3m+k\}\\
 &\cup\{3m+2k+1,\ldots,L\}.
\end{split}
\end{equation}
Their signs are $+,-,+$. Write $c_0=(k+1)/2$ and, in observation order, define
\begin{align}
 z^+_{i}&=r(c_0+m-i),&&1\le i\le m,\nonumber\\
 z^-_{i}&=-r[c_0+\min(i-1,n-i)],&&1\le i\le n,\nonumber\\
 z^{+,R}_{i}&=r(c_0+i-1),&&1\le i\le m.\label{eq:localz}
\end{align}
Let $a\ge r(k+n-1)/2$, $s_i^0=a\epsilon_i$, and $w=s^0-z$. Within a half-window the slope magnitude is $r$; across a transition the change is $2rc_0=r(k+1)$; the central edge of the negative window is flat. Thus $z$ is rate-feasible. Each level appears twice with each sign, so $\sum_iw_i=0$.

Let $Q_i=\sum_{j\le i}w_j$. The left-side partial sums are nonnegative, reach zero at the center of the negative window, and become nonpositive on the reflected right side. Whenever $Q_i\ne0$, the corresponding slope of $z$ equals $-r\operatorname{sign}(Q_i)$ per unit time. Summation by parts gives
\begin{equation}
 h_{\K_I(\infty,r)}(w)
 =r\sum_i(t_{i+1}-t_i)|Q_i|=w^Tz.
\end{equation}
Equivalently, $q_i=-Q_i$ satisfies $z-s^0+D^Tq=0$ and all edge complementarity conditions. The convex KKT conditions show that $z$ is the unique local constant-target projection. Consequently $2w^Ts^0-\norm w^2-2h_{\K_I}(w)=\norm w^2$, and direct summation proves \eqref{eq:Lc}.

If the block begins at slot $x+1$, its actual target is
\begin{equation}
 s_i=\epsilon_i[a+\eta(t_i-c)],\qquad a=\rho_0+\eta(x+c).
\end{equation}
Both $w$ and $\epsilon$ are reflection-symmetric, while $t_i-c$ is antisymmetric. Hence
\begin{equation}
 \sum_iw_i\epsilon_i(t_i-c)=0,
 \quad w^Ts=w^Ts^0,
 \quad\norm s^2=2na^2+\eta^2M_2,
\end{equation}
where $M_2=\sum_{i\in I}(t_i-c)^2=O((n+k)^3)$. This is an exact affine cancellation; the same vector is a feasible dual certificate for the affine target.

Concatenate the local dual directions, assigning zero weight to the healthy prefix, any early blocks that fail the amplitude condition, and the final unused reads. Every feasible whole-history path restricts to a feasible path on each block. Therefore \eqref{eq:support} holds even if the particular local primal optimizers cannot be joined. The elementary weak dual inequality
\begin{equation}
 \min_{z\in\K}\norm{s-z}^2
 \ge2w^Ts-\norm w^2-2h_\K(w)
\end{equation}
then proves \eqref{eq:closedG} and
\begin{equation}
2\sigma^2D_H\le E_{\rm dead}+E_{\rm tail}
+\sum_\ell[L_c(a_\ell,n_\ell)+\eta^2M_{2,\ell}].\label{eq:sharpfinite}
\end{equation}
This concatenates dual directions and uses restriction of the global feasible class.

For the terminal-balanced construction, $M=\sqrt{H/\xi}+O(1)$, $n_\ell=\xi\ell+O(1)$, and $a_\ell=\eta\xi\ell^2+O(\ell+1)$. Each block is perturbed by $O(1)$ readings, and fewer than four tail slots remain. The leading sums are
\begin{align}
 E_{\rm dead}&=\frac{2k\eta^2}{5\sqrt\xi}H^{5/2}+O(H^2),\nonumber\\
 \sum_\ell L_c(a_\ell,n_\ell)&=\frac{r\eta\sqrt\xi}{5}H^{5/2}+O(H^2),\nonumber\\
 \eta^2\sum_\ell M_{2,\ell}+E_{\rm tail}&=O(H^2).
\end{align}
Indeed, the first two per-block main terms are $2k\eta^2\xi^2\ell^4$ and $r\eta\xi^3\ell^4$; all other terms sum to $O(M^4)$. The finite number of early zero-weight blocks changes only the constant. Substitution in \eqref{eq:sharpfinite} proves \eqref{eq:upper}. Its coefficient is minimized at $\xi=2k\eta/r$, where it equals $2\sqrt2\sqrt{kr}\eta^{3/2}/5$. Combining with \eqref{eq:sharp_lower} proves the exact limit.
\end{proof}


\section{Proof of the Finite-Box Theorem}
\begin{proof}
Write $b_j=2Ba_j-B^2$ and $N_B=\lceil2B/r\rceil+1$. Choose a healthy explanation equal to zero on the prefix. After onset it moves at speed at most $r$ toward the endpoint $\epsilon B$ of the currently observed task, and holds there on arrival. At a task change it starts moving toward the new endpoint. It never has to return to zero at the final time. Interpolation during excluded slots preserves the same rate bound.

For an observed point with target sign $\epsilon$, the deficiency in squared-energy saving relative to $b_j$ is
\begin{equation}
 b_j-(2\epsilon a_jz_j-z_j^2)
 =2a_j(B-\epsilon z_j)-(B^2-z_j^2)\le4Ba_j.
\end{equation}
It is zero once the tracker reaches $\epsilon B$. Following a task change, at most $N_B$ observations can be deficient before the tracker reaches that endpoint; an intervening change restarts the count. Assign each deficient observation to its latest change. Every internal change can be paired with a disjoint complete $k$-slot movement immediately preceding the new observed run, starting at $m$. Its assigned deficiency is at most
\begin{equation}
 4BN_B\{a_m+\eta(k+N_B)\}.
\end{equation}
On an excluded slot the oracle energy exceeds the baseline by $(a_j-B)^2$. The paired move thus supplies at least $k(a_m-B)_+^2$; unpaired excluded slots supply nonnegative excess. A change can therefore have net adverse contribution at most
\begin{equation}
 q_+(a_m)=\pos{4BN_B[a_m+\eta(k+N_B)]-k(a_m-B)_+^2}.
\end{equation}
Because $\eta>0$, $q_+(a_m)$ is nonzero at only finitely many possible start indices. Define the unwhitened constant
\begin{equation}
 C_1=4BN_B(\rho_0+\eta N_B)+\sum_{m\ge1}q_+(\rho_0+\eta m).
\end{equation}
The sum has finitely many nonzero terms. The initial approach from zero contributes the first term, and the tracker stays at its last endpoint. Simultaneously for all calendars,
\begin{equation}
 2\sigma^2D_H^B(\pi)\ge\sum_{j=1}^Hb_j-C_1.\label{eq:finiteLower}
\end{equation}
This is a uniform full-history construction.

For staying, dropping the rate constraints and the nonnegative prefix loss gives
\begin{equation}
 2\sigma^2G_H^B(\mathrm{stay})\ge\sum_{j=1}^H(a_j-B)_+^2.
\end{equation}
Consequently its deficit is at most $\sum b_j/(2\sigma^2)+C_2$, where
\begin{equation}
 2\sigma^2C_2=\sum_{j\ge1:a_j<B}(B-a_j)^2<\infty.
\end{equation}
Conversely, choosing the feasible constant $z=B$ throughout the prefix and subsequent stay experiment gives
\begin{equation}
2\sigma^2D_H^B(\mathrm{stay})\ge\sum_{j=1}^Hb_j-n_0B^2.
\end{equation}
The two inequalities identify its full polynomial baseline up to $O(1)$. Combining the stay upper bound with \eqref{eq:finiteLower}, and using optimality relative to staying, proves both \eqref{eq:finite} and \eqref{eq:regret}.
\end{proof}

An explicit bound in the theorem is
\begin{equation}
 C_B=\frac{C_1+\sum_{j\ge1:a_j<B}(B-a_j)^2}{2\sigma^2}.
 \label{eq:explicit_CB}
\end{equation}
Both terms in the numerator are unwhitened squared-torque quantities. For the fast-ramp example at $\sigma=.2$ this gives about $2.97\times10^5$ nats, compared with a $14.7$-nat advantage in the tested switching family. It is a uniform existence bound, while finite-horizon designs are evaluated by the exact projection.

\section{Full-Vector Rank-One Reduction}\label{app:rankone}
\begin{proposition}[One shared healthy direction]\label{prop:rankone}
Assume the healthy loading is exactly $+h$ and $-h$ in the two tasks, and the complete independent Gaussian observations have the same known covariance $\Sigma\succ0$ in both tasks. Consider
$Y_i=\epsilon_i h\beta_i+fa_i+\xi_i$, and the same deterministic calendar class as Theorem~\ref{thm:sharp}. Let $g=\Sigma^{-1/2}h\ne0$, $v=\Sigma^{-1/2}f$, $c_f=g^Tv/\|g\|$, $F=\|v\|^2$, and $r_g>0$ be the difference-path rate of $\|g\|\beta$ per slot. If $F>0$ and $c_f\ne0$, the unbounded-level optimal deficit satisfies
\begin{equation}
 D_H^*\sim\frac{\sqrt2}{5}\sqrt{kF r_g|c_f|}\,\eta^{3/2}H^{5/2}.\label{eq:rankone}
\end{equation}
A matching planned construction has $\xi=2kF\eta/(r_g|c_f|)$. If $c_f=0$, staying observes all unconfounded fault energy and has zero deficit.
\end{proposition}
\begin{proof}
Multiply each complete observation by its known task sign, whiten, and split orthogonally along $g/\|g\|$. With $I_\pi$ the retained slots, this gives the exact identity
\begin{align}
G_H(\pi)={}&\frac{F-c_f^2}{2}\sum_{i\in I_\pi}a_i^2\nonumber\\
&+\frac12\min_{z\in\K_\pi(\infty,r_g)}\sum_{i\in I_\pi}(c_f\epsilon_i a_i-z_i)^2.\label{eq:vector_split}
\end{align}
The healthy prefix remains in the scalar projection with zero target. Hence excluded slots cost $Fa_j^2/2$, whereas only the component $c_f$ can be explained by the healthy path. Changing its sign preserves the symmetric path class; assume $c_f>0$.

For the converse, repeat the late-block construction in Appendix~B with minimum amplitude $A$. Use $z_i=\epsilon_i r_g\min(i,n-1-i)$ on each clipped observed run, and zero at all artificial boundaries. Put $C=r_g(2c_f A-r_gL/2)/4$. For $C\ge0$, twice the deficit on that block is bounded below by
\begin{equation}
 FA^2d+\frac{Ck(L-d)^2}{d+k}-2CL.
\end{equation}
AM--GM has leading term $\sqrt{2kF r_gc_f}\,A^{3/2}L$. The same $L=\lfloor H^{3/4}\rfloor$ partition makes all boundary and amplitude-variation terms lower order uniformly in the calendar. This yields the liminf coefficient $\sqrt2\sqrt{kF r_gc_f}\eta^{3/2}/5$ in \eqref{eq:rankone}.

For the upper bound, use the symmetric local scalar witness of Appendix~B with amplitude $c_f a$, slope $c_f\eta$, and rate $r_g$. In the orthogonal subspace use the full matched direction $\epsilon_i a_i(v-c_fg/\|g\|)$, whose healthy support is zero. Their concatenation gives
\begin{equation}
2D_H\le\frac15\left(\frac{2kF\eta^2}{\sqrt\xi}+r_gc_f\eta\sqrt\xi\right)H^{5/2}+O(H^2).
\end{equation}
Minimizing this coefficient at $\xi=2kF\eta/(r_gc_f)$ matches the converse. If $c_f=0$, \eqref{eq:vector_split} has no fault in the scalar subproblem, and observing every slot attains the full-calendar oracle. The case $g=0$ is likewise unconfounded. The argument uses one shared healthy direction.
\end{proof}

\section{Extended Operating-Range Certificate}\label{app:embedding}
This appendix verifies the stated synthetic rigid model without changing its controller or transition duration. Let $\omega=40$, $b_m=.02$, $r_m=.001$, $f_{\max}=.75$ and $\dot f_{\max}=.003$ in the units of Section~\ref{sec:control}. The mass and bias convention is exactly \eqref{eq:plant}. Signals $\beta_m,f$ need only be absolutely continuous; their first derivatives are bounded almost everywhere, with no curvature assumption.

\subsection{Global geometry and the moving reference}
Using segment indices $s=1,\ldots,6$, for each link center or payload point let $\ell_s^{(j)}$ be its contributing segment lengths, $R_i^{(j)}=\sum_{s\ge i}\ell_s^{(j)}$, and
\begin{equation}
 L_{J,j}^2=\sum_i\left(\sum_{s\ge i}\ell_s^{(j)}\sqrt{s}\right)^2.
\end{equation}
The Cartesian Jacobian satisfies $\|J_j\|\le\|R^{(j)}\|$ and $\|D J_j(q)v\|_F\le L_{J,j}\|v\|$. Consequently $L_M=\sum_j2m_j\|R^{(j)}\|L_{J,j}$ bounds the Lipschitz constant of $M_0$. The constant rotor and link rotational terms add no derivatives. Taking row sums of the entrywise matrix $9.81\sum_{s\ge\max(i,j)}\ell_s$ bounds the payload gravity derivative. Cumulative angular speeds obey $(\sum_{i\le s}\dot q_i)^2\le s\|\dot q\|^2$. These formulas give conservative bounds
\begin{equation}
\begin{gathered}
 d_0=.5,\quad \bar A=1.2908,\quad G=11.145473,\\
 L_M=4.489123,\quad L_g=23.3478,\quad C_p=2.703999,\\
 \bar M=2.389402,\quad C_c=2.572884.
\end{gathered}\label{eq:geometry_caps}
\end{equation}
Here $M_0,M\succeq d_0I$, $\|A\|\le\bar A$, $\|g_p\|\le G$, $\|c_p\|\le C_p\|\dot q\|^2$, $\|M\|\le\bar M$, and $\|c(q,\dot q,m)\|\le C_c\|\dot q\|^2$.

For $x=(\omega e,\dot e)$ use $P=\left[\begin{smallmatrix}1.5&.5\\.5&.5\end{smallmatrix}\right]$ and $A_0=\left[\begin{smallmatrix}0&1\\-1&-2\end{smallmatrix}\right]$, so $A_0^TP+PA_0=-I$. With $\mu=b_m\bar A/d_0$, the exact tracking dynamics and $\|M^{-1}\beta_m A\|\le\mu$ imply
\begin{equation}
 \tfrac{\mathrm d}{\mathrm dt}\sqrt{x^T(P\otimes I)x}
 \le-\lambda_0\sqrt{x^T(P\otimes I)x}
       +\frac{\|Pb_2\|}{\sqrt{p_-}}\bar d,
\end{equation}
where $b_2=(0,1)^T$, $p_-=\lambda_{\min}(P)$, $p_+=\lambda_{\max}(P)$, and
\begin{equation}
 \lambda_0(\omega)=\frac{\omega\{1-2\|Pb_2\|\sqrt5\,\mu\}}{2p_+}.
 \label{eq:lambda_control}
\end{equation}
Indeed, the state-dependent inertia perturbation is at most $\mu\omega\sqrt5\|x\|$, since $\|\omega^2e+2\omega\dot e\|\le\omega\sqrt5\|x\|$. Its contribution is absorbed into the Lyapunov decay. Here $\|Pb_2\|=1/\sqrt2$, $p_+=1+1/\sqrt2$, so $\lambda_0(40)>9.8028$ and the rounded lower rate $9.8$ is valid. On $\|x\|\le X=.6$,
\begin{equation}
 \bar d=\frac{b_m(\bar A a_d+C_p(v_d+X)^2+G)+r_m\bar A(v_d+X)+f_{\max}}{d_0}.
\end{equation}
The quintic reference gives $v_d\le1.206414$, $a_d\le2.476528$. From $x(0)=0$ the resulting upper bound is $.598945<X$, closing a global bootstrap across every predetermined transition. Let $\mathcal M_0$ bound $|M_0|$ entrywise by the same reach-vector outer products, and let $C_{0,i}$ and $\mathcal G_{0,i}$ bound the nominal centrifugal and gravity components. Then
\begin{equation}
 |\tau_{u,i}|\le\|\mathcal M_{0,i:}\|_2(a_d+\omega\sqrt5X)
 +C_{0,i}(v_d+X)^2+\mathcal G_{0,i}.
\end{equation}
This gives $(96.747,70.165,50.819,38.440,31.864,29.049)$ N m, below the specified limits. Thus the unsaturated calculation is self-consistent.

\subsection{Equilibrium shift and rate-driven comparison}
During a hold define
\begin{equation}
 d(q,\beta_m,f)=-M_0(q)^{-1}\{\beta_m g_p(q)+e_f f\}.
\end{equation}
Put $D_0=(b_mG+f_{\max})/d_0$, $D_1=(r_mG+\dot f_{\max})/d_0$ and
$L_d=L_M(b_mG+f_{\max})/d_0^2+b_mL_g/d_0$. Since $L_d<\omega^2$, $q_e=q_d+d(q_e)/\omega^2$ is a contraction on the radius-$D_0/\omega^2$ ball. It has a unique solution, with
\begin{equation}
 \|q_e-q_d\|\le D_0/\omega^2,\qquad
 \|\dot q_e\|\le D_1/(\omega^2-L_d).
\end{equation}
The exact hold dynamics are
\begin{align}
 \ddot q+2\omega\dot q+\omega^2(q-q_d)&=d+z,\nonumber\\
 z=-M_0^{-1}(\beta_m A\ddot q&+\beta_m c_p+\dot\beta_m A\dot q).
\end{align}
Thus $\|z\|\le\mu\|\ddot q\|+d_v\|\dot q\|$, with
$d_v=(b_mC_pX+r_m\bar A)/d_0$, and
$\|\dot d\|\le D_1+L_d\|\dot q\|$. No derivative of $z$ is used.

For $0\le\ell<\omega$, the exponentially weighted absolute gains of the critical-damping kernels are
\begin{align}
 g_e(\ell)&=(\omega-\ell)^{-2},\nonumber\\
 g_v(\ell)&=\frac{-\ell+2\omega e^{-(\omega-\ell)/\omega}}{(\omega-\ell)^2},\nonumber\\
 g_a(\ell)&=1+\frac{2\omega}{\omega-\ell}
 +\frac{\omega^2(2e^{-2(\omega-\ell)/\omega}-1)}{(\omega-\ell)^2}.
\end{align}
The $1$ in $g_a$ retains the zero-time atom in the acceleration response. Define
\begin{equation}
 K_\ell=\begin{bmatrix}
 \mu g_a&L_dg_v+d_vg_a\\
 \mu g_v&L_dg_e+d_vg_v
 \end{bmatrix}.
\end{equation}
For the homogeneous response, $P_0=X/\omega+D_0/\omega^2$ and $V_0=X$ bound the initial shifted position and speed. Upper bounds after multiplication by $e^{\ell t}$ are
\begin{align}
 h_A&=\omega^2P_0+2\omega V_0+\frac{\omega^2(V_0+\omega P_0)}{e(\omega-\ell)},\nonumber\\
 h_V&=V_0+\frac{\omega(V_0+\omega P_0)}{e(\omega-\ell)}.
\end{align}
Direct outward arithmetic verifies the positive supersolutions
\begin{align}
 (I-K_0)\begin{bmatrix}.00061\\.000019\end{bmatrix}
 &\ge D_1\begin{bmatrix}g_v(0)\\g_e(0)\end{bmatrix},\nonumber\\
 (I-K_{12})\begin{bmatrix}130\\1.7\end{bmatrix}
 &\ge\begin{bmatrix}h_A\\h_V\end{bmatrix}.
\end{align}
The componentwise relative margins in the first supersolution inequality are about $1.68\%$ and $2.81\%$. These are arithmetic feasibility margins, not additional model-uncertainty allowances. Both nonnegative matrices have spectral radius below one. Positive Volterra comparison, including the acceleration feedthrough, gives
\begin{equation}
 \|\ddot q(t)\|\le130e^{-12t}+.00061,\quad
 \|\dot q(t)\|\le1.7e^{-12t}+.000019.\label{eq:rate_tube}
\end{equation}
This holds from every hold start under the global transition tube. The remaining position estimate follows directly from the equilibrium equation:
\begin{equation}
 \|q-q_e\|\le
 \frac{(1+\mu)\|\ddot q\|+(2\omega+d_v)\|\dot q\|}{\omega^2-L_d}.
\end{equation}
At $t\ge1.5$ s this gives the bounds in Section~\ref{sec:control}. Finally, the residual error relative to the reference-pose static model is bounded by
\begin{align}
 \|e_{\rm red}\|\le{}&\bar M\|\ddot q\|+C_c\|\dot q\|^2\nonumber\\
 &+r_m\bar A\|\dot q\|+b_mL_g\|q-q_d\|\nonumber\\
 <{}&.002031\ {\rm N\,m}.
\end{align}
All static amplitude dependence remains in $D_0,L_d$ and the effort budget; it has not been discarded. The small persistent dynamic error is driven by bounded rates rather than by treating a static torque as arbitrary time-varying forcing. For completeness, write $A_0=.00061$, $A_1=130$, $V_0=.000019$, $V_1=1.7$, and $L=\omega^2-L_d$. Define
\begin{align}
 Q_0&=D_0/\omega^2+\{(1+\mu)A_0+(2\omega+d_v)V_0\}/L,\nonumber\\
 Q_1&=\{(1+\mu)A_1+(2\omega+d_v)V_1\}/L.\nonumber
\end{align}
Substitution in the residual bound gives \eqref{eq:admission_constants} with
\begin{align}
 E_\infty&=\bar M A_0+C_cV_0^2+r_m\bar A V_0+b_mL_gQ_0,\nonumber\\
 C_1&=\bar M A_1+2C_cV_0V_1+r_m\bar A V_1+b_mL_gQ_1,\nonumber\\
 C_2&=C_cV_1^2.\nonumber
\end{align}
Since $e^{-24t}\le e^{-12t}$ for $t\ge0$, $C=C_1+C_2$ is a valid single-exponential coefficient.

\begin{thebibliography}{99}
\bibitem{chernoff} H. Chernoff, ``Sequential design of experiments,'' \emph{Ann. Math. Statist.}, vol. 30, no. 3, pp. 755--770, 1959.
\bibitem{naghshvar} M. Naghshvar and T. Javidi, ``Active sequential hypothesis testing,'' \emph{Ann. Statist.}, vol. 41, no. 6, pp. 2703--2738, 2013, doi: 10.1214/13-AOS1144.
\bibitem{atkinson} A. C. Atkinson and V. V. Fedorov, ``The design of experiments for discriminating between two rival models,'' \emph{Biometrika}, vol. 62, no. 1, pp. 57--70, 1975, doi: 10.1093/biomet/62.1.57.
\bibitem{burnashev} M. V. Burnashev, ``On the minimax detection of an inaccurately known signal in a white Gaussian noise background,'' \emph{Theory Probab. Appl.}, vol. 24, no. 1, pp. 107--119, 1979, doi: 10.1137/1124008.
\bibitem{gjn} A. Goldenshluger, A. Juditsky, and A. Nemirovski, ``Hypothesis testing by convex optimization,'' \emph{Electron. J. Statist.}, vol. 9, no. 2, pp. 1645--1712, 2015, doi: 10.1214/15-EJS1054.
\bibitem{basseville} M. Basseville and I. V. Nikiforov, \emph{Detection of Abrupt Changes: Theory and Application}. Englewood Cliffs, NJ, USA: Prentice Hall, 1993.
\bibitem{zhang2002} X. Zhang, M. M. Polycarpou, and T. Parisini, ``A robust detection and isolation scheme for abrupt and incipient faults in nonlinear systems,'' \emph{IEEE Trans. Autom. Control}, vol. 47, no. 4, pp. 576--593, 2002, doi: 10.1109/9.995036.
\bibitem{nik98} R. Nikoukhah, ``Guaranteed active failure detection and isolation for linear dynamical systems,'' \emph{Automatica}, vol. 34, no. 11, pp. 1345--1358, 1998.
\bibitem{campbell} S. L. Campbell and R. Nikoukhah, \emph{Auxiliary Signal Design for Failure Detection}. Princeton, NJ, USA: Princeton University Press, 2004.
\bibitem{scott} J. K. Scott, R. Findeisen, R. D. Braatz, and D. M. Raimondo, ``Input design for guaranteed fault diagnosis using zonotopes,'' \emph{Automatica}, vol. 50, no. 6, pp. 1580--1589, 2014, doi: 10.1016/j.automatica.2014.03.016.
\bibitem{blanchini} F. Blanchini, D. Casagrande, G. Giordano, S. Miani, S. Olaru, and V. Reppa, ``Active fault isolation: A duality-based approach via convex programming,'' \emph{SIAM J. Control Optim.}, vol. 55, no. 3, pp. 1619--1640, 2017, doi: 10.1137/15M1046046.
\bibitem{heirung} T. A. N. Heirung and A. Mesbah, ``Input design for active fault diagnosis,'' \emph{Annu. Rev. Control}, vol. 47, pp. 35--50, 2019, doi: 10.1016/j.arcontrol.2019.03.002.
\bibitem{simandl} M. \v{S}imandl and I. Pun\v{c}och\'{a}\v{r}, ``Active fault detection and control: Unified formulation and optimal design,'' \emph{Automatica}, vol. 45, pp. 2052--2059, 2009.
\bibitem{ashari} A. Esna Ashari, R. Nikoukhah, and S. L. Campbell, ``Effects of feedback on active fault detection,'' \emph{Automatica}, vol. 48, pp. 866--872, 2012.
\bibitem{coster} D. C. Coster and C.-S. Cheng, ``Minimum cost trend-free run orders of fractional factorial designs,'' \emph{Ann. Statist.}, vol. 16, no. 3, pp. 1188--1205, 1988.
\bibitem{cheng} C.-S. Cheng and D. M. Steinberg, ``Trend robust two-level factorial designs,'' \emph{Biometrika}, vol. 78, no. 2, pp. 325--336, 1991.
\bibitem{bickel1} P. J. Bickel and A. M. Herzberg, ``Robustness of design against autocorrelation in time I: Asymptotic theory, optimality for location and linear regression,'' \emph{Ann. Statist.}, vol. 7, no. 1, pp. 77--95, 1979.
\bibitem{bickel2} P. J. Bickel, A. M. Herzberg, and M. F. Schilling, ``Robustness of design against autocorrelation in time II: Optimality, theoretical and numerical results for the first-order autoregressive process,'' \emph{J. Amer. Statist. Assoc.}, vol. 76, no. 376, pp. 870--877, 1981.
\bibitem{schieder} R. Schieder and C. Kramer, ``Optimization of heterodyne observations using Allan variance measurements,'' \emph{Astron. Astrophys.}, vol. 373, pp. 746--756, 2001.
\bibitem{ossenkopf} V. Ossenkopf, ``The stability of spectroscopic instruments: A unified Allan variance computation scheme,'' \emph{Astron. Astrophys.}, vol. 479, pp. 915--926, 2008.
\bibitem{vitus} M. P. Vitus, W. Zhang, A. Abate, J. Hu, and C. J. Tomlin, ``On efficient sensor scheduling for linear dynamical systems,'' \emph{Automatica}, vol. 48, no. 10, pp. 2482--2493, 2012.
\bibitem{zhao} L. Zhao, W. Zhang, J. Hu, A. Abate, and C. J. Tomlin, ``On the optimal solutions of the infinite-horizon linear sensor scheduling problem,'' \emph{IEEE Trans. Autom. Control}, vol. 59, no. 10, pp. 2825--2830, 2014.
\bibitem{banerjee} T. Banerjee and V. V. Veeravalli, ``Data-efficient quickest change detection with on--off observation control,'' \emph{Sequential Analysis}, vol. 31, no. 1, pp. 40--77, 2012.
\bibitem{lautay} T. S. Lau and W. P. Tay, ``Asymptotically optimal sampling policy for quickest change detection with observation-switching cost,'' \emph{IEEE Trans. Signal Process.}, vol. 69, pp. 1332--1346, 2021, doi: 10.1109/TSP.2021.3057258.
\bibitem{lubenia} P. V. N. Lubenia and T. Banerjee, ``Quickest change detection with cost-constrained experiment design,'' arXiv:2509.14186v2, 2025, preprint.
\bibitem{liang} Y. Liang, A. G. Tartakovsky, and V. V. Veeravalli, ``Quickest change detection with non-stationary post-change observations,'' \emph{IEEE Trans. Inf. Theory}, vol. 69, no. 5, pp. 3400--3414, 2023, doi: 10.1109/TIT.2022.3230583.
\bibitem{hou} Y. Hou, Y. Oleyaeimotlagh, R. Mishra, H. Bidkhori, and T. Banerjee, ``Robust quickest change detection in nonstationary processes,'' \emph{Sequential Analysis}, vol. 43, no. 3, pp. 275--300, 2024, doi: 10.1080/07474946.2024.2356555.
\bibitem{deluca} A. De Luca and R. Mattone, ``Actuator failure detection and isolation using generalized momenta,'' in \emph{Proc. IEEE Int. Conf. Robot. Autom.}, vol. 1, 2003, pp. 634--639, doi: 10.1109/ROBOT.2003.1241665.
\bibitem{haddadin} S. Haddadin, A. De Luca, and A. Albu-Sch\"affer, ``Robot collisions: A survey on detection, isolation, and identification,'' \emph{IEEE Trans. Robot.}, vol. 33, no. 6, pp. 1292--1312, 2017, doi: 10.1109/TRO.2017.2723903.
\bibitem{lynch} K. M. Lynch and F. C. Park, \emph{Modern Robotics: Mechanics, Planning, and Control}. Cambridge, U.K.: Cambridge University Press, 2017, chs. 8 and 11.
\end{thebibliography}

\end{document}
